\cotefile{157-3}
\scanpage{192}\mspage{141}
et les flèches correspondantes

\diagram{\(\left\{\begin{aligned}
\varphi_Q &: \unsure{\mathscr X|\mathcal X}_{\deleted{\unsure{Q}}Q}=\left(\int \begin{tikzcd}[ampersand replacement=\&,row sep=0pt,column sep=1.2em,cramped] \& X_1 \\ X_0\arrow[ur,dash]\arrow[dr] \& \\ \& X_2 \end{tikzcd}\right)\longrightarrow Y\\
\varphi_{Q^\circ} &: \mathscr X^{\deleted{\unsure{Q}}}_{Q^\circ}=\left(\int \begin{tikzcd}[ampersand replacement=\&,row sep=0pt,column sep=1.2em,cramped] \& X_1^\circ \\ X_0^\circ\arrow[ur,dash]\arrow[dr,dash] \& \\ \& X_2^\circ \end{tikzcd}\right)\longrightarrow Y^\circ .
\end{aligned}\right.\)}

Enfin, considérons aussi les flèches $\widetilde\varphi_Q,\inserted{\unsure{\overline\varphi_{Q^\circ}|\widetilde\varphi_{Q^\circ}}}$ canoniques, relatives à l'intégrale propre $\int_{\mathrm{pr}}$ (au lieu de l'intégrale lisse $\int=\int_{\unsure{\mathrm{lis}|\mathrm{lic}}}$)

\diagram{\(\left\{\begin{aligned}
\widetilde\varphi_Q &\unsure{\mathrel{:}|\mathrel{=}} \widetilde{\mathscr X}_Q=\left(\int_{\mathrm{pr}} \begin{tikzcd}[ampersand replacement=\&,row sep=0pt,column sep=1.2em,cramped] \& X_1 \\ X_0\arrow[ur,dash]\arrow[dr] \& \\ \& X_2 \end{tikzcd}\right)\longrightarrow Y\\
\deleted{\text{\unsure{associées}}}\quad\deleted{\widetilde\varphi_{Q^\circ}:}&\deleted{\left(\int_{\mathrm{pr}} \begin{tikzcd}[ampersand replacement=\&,row sep=0pt,column sep=1.2em,cramped] \& X_1^\circ \\ X_0^\circ\arrow[ur,dash]\arrow[dr,dash] \& \\ \& X_2^\circ \end{tikzcd}\right)\longrightarrow \unsure{Y^\circ}}\\
\widetilde\varphi_{Q^\circ} &: \widetilde{\mathscr X}_{Q^\circ}=\left(\int \begin{tikzcd}[ampersand replacement=\&,row sep=0pt,column sep=1.2em,cramped] \& X_1^\circ \\ X_0^\circ\arrow[ur]\arrow[dr,dash] \& \\ \& X_2^\circ \end{tikzcd}\right)\longrightarrow Y^\circ .
\end{aligned}\right.\)}

On \unsure{a} alors ce qui suit :

\[
\text{\textcircled{\unsure{a}}}\quad \varphi_{Q^\circ}\simeq(\widetilde\varphi_Q)^\circ,\quad \varphi_Q\simeq(\widetilde\varphi_{Q^\circ})^\circ
\]

(d'où résulte déjà que $\left\{\begin{aligned}\varphi_Q\in W&\Longleftrightarrow\widetilde\varphi_{Q^\circ}\in W\\\varphi_{Q^\circ}\in W&\Longleftrightarrow\widetilde\varphi_Q\in W\end{aligned}\right.$ \unsure{\deleted{)}} ${}^{\star}$
% slanted
\margin{et que les flèches $\varphi_Q$ et $\widetilde\varphi_{Q^\circ}$ d'une part, $\varphi_{Q^\circ}$ et $\widetilde\varphi_Q$ d'autre part, sont \unsure{isomorphes} comme flèches dans $\mathrm{Hot}_W$.}).

\textcircled{b} Les relations

\[
\varphi_Q\in W,\quad \varphi_{Q^\circ}\in W,\quad \widetilde\varphi_Q\in W,\quad \widetilde\varphi_{Q^\circ}\in W
\]

sont équivalentes, en d'autres termes : pour que le carré soit $W$-homotopiquement cocartésien (version intégrale lisse), il f. et il suffit que le carré opposé $Q^\circ$ le soit, et ces conditions équivalent aussi aux conditions correspondantes en termes des intégrales propres (au lieu des intégrales lisses).

\textcircled{c} Les quatre flèches $\varphi_Q,\varphi_{Q^\circ},\widetilde\varphi_Q,\widetilde\varphi_{Q^\circ}$, \inserted{\unsure{de $\mathrm{Cat}$|de $\mathrm{Cat}$,} en tant que flèches} dans $\mathrm{Hot}_W$, sont isomorphes dans \uline{$\mathrm{Fl}$}$(\mathrm{Hot}_W)$. Plus précisément, en tant que flèches de $\mathrm{Cat}$,
% archive sheet number: 192
